\documentclass[11pt,a4paper]{article}
\usepackage[utf8]{inputenc}
\usepackage{amsmath}
\usepackage{amsfonts}
\usepackage{amssymb}
\usepackage{geometry}

\geometry{top=2.5cm, bottom=2.5cm, left=2.5cm, right=2.5cm}

\title{\textbf{Penyelesaian Panjang Busur Kurva Polar}}
\author{Ahmad Fakhrul Bawani}
\date{}

\begin{document}

\maketitle

\section{Deskripsi Masalah}
Diberikan sebuah persamaan kurva polar:
\begin{equation}
  r = \sqrt{1 + \cos(2\theta)}
\end{equation}
Hitunglah panjang lintasan (\textit{arc length}) dari kurva tersebut pada rentang sudut:
\begin{equation*}
  0 \le \theta \le \frac{\pi\sqrt{2}}{4}
\end{equation*}

\section{Penyederhanakan Fungsi Kurva}
Menggunakan identitas sudut ganda trigonometri $\quad 1 + \cos(2\theta) = 2\cos^2\theta \quad$, persamaan kurva polar dapat disederhanakan terlebih dahulu:
\begin{equation*}
  r = \sqrt{2\cos^2\theta} = \sqrt{2}|\cos\theta|
\end{equation*}
Karena nilai batas atas sudut $\theta = \frac{\pi\sqrt{2}}{4} \approx 1.11\text{ rad}$ berada di dalam Kuadran I ($0 \le \theta \le \frac{\pi}{2}$), maka fungsi $\cos\theta$ bernilai positif. Oleh karena itu, tanda mutlak dapat dihilangkan:
\begin{equation}
  r = \sqrt{2}\cos\theta
\end{equation}

\section{Langkah-Langkah Perhitungan}

\subsection{Menentukan Turunan dan Komponen Kuadrat}
Turunan pertama fungsi $r$ terhadap $\theta$ adalah:
\begin{equation*}
  \frac{dr}{d\theta} = -\sqrt{2}\sin\theta
\end{equation*}

Kita kuadratkan masing-masing komponen untuk dimasukkan ke dalam bentuk radikan panjang busur:
\begin{align*}
  r^2 &= (\sqrt{2}\cos\theta)^2 = 2\cos^2\theta \\
  \left(\frac{dr}{d\theta}\right)^2 &= (-\sqrt{2}\sin\theta)^2 = 2\sin^2\theta
\end{align*}

\subsection{Menyederhanakan Bentuk di Dalam Akar}
Gabungkan kedua komponen kuadrat tersebut dengan menggunakan identitas dasar Pythagoras trigonometri $\cos^2\theta + \sin^2\theta = 1$:
\begin{align*}
  r^2 + \left(\frac{dr}{d\theta}\right)^2 &= 2\cos^2\theta + 2\sin^2\theta \\
  &= 2(\cos^2\theta + \sin^2\theta) = 2
\end{align*}

\subsection{Evaluasi Integral Panjang Busur}
Rumus panjang busur dalam koordinat polar didefinisikan sebagai:
\begin{equation}
  L = \int_{\alpha}^{\beta} \sqrt{r^2 + \left(\frac{dr}{d\theta}\right)^2} \, d\theta
\end{equation}

Substitusikan batas integral $\alpha = 0$ dan $\beta = \frac{\pi\sqrt{2}}{4}$ serta hasil penyederhanaan radikan yang telah diperoleh:
\begin{align*}
  L &= \int_{0}^{\frac{\pi\sqrt{2}}{4}} \sqrt{2} \, d\theta \\
  L &= \sqrt{2} \ \Big[ \theta \Big]_{0}^{\frac{\pi\sqrt{2}}{4}} \\
  L &= \sqrt{2} \left( \frac{\pi\sqrt{2}}{4} - 0 \right) \\
  L &= \frac{2\pi}{4} = \frac{\pi}{2}
\end{align*}

\section{Kesimpulan}
Panjang busur total dari kurva $r = \sqrt{1 + \cos(2\theta)}$ pada rentang sudut yang ditentukan adalah tepat \textbf{$\frac{\pi}{2}$} satuan panjang.

\end{document}